% Standalone copy of the rewritten Introduction of the manuscript
% FINAL_SUBMISSION.tex (Section 1). The TikZ workflow figure that follows
% the last paragraph in the manuscript is not repeated here; every label,
% citation key and cross-reference is as in the manuscript.

\section{Introduction}\label{sec:intro}

The hull of a linear code $C$ is the intersection $C\cap C^{\perp}$ with its
dual, and its dimension is a structural invariant of the code. Replacing the
Euclidean inner product by a Frobenius-twisted pairing in one argument leads to
the $k$-Galois inner product and to the $k$-Galois hull, of which the Euclidean
and Hermitian hulls are the cases $k=0$ and $k$ equal to half the degree of the
field \cite{fanZhang2017}. For cyclic-type codes the ambient space is a
polynomial quotient, so duals, intersections and hence hull dimensions are
accessible to factorization arguments \cite{huffmanPless2003}. Hull dimensions
therefore describe a single code and, at the same time, define a statistic of an
entire family; this paper studies the joint distribution of that statistic and
the code dimension.

The hull dimension, and not only its vanishing, controls several constructions.
The algorithms that test permutation equivalence of two linear codes, and those
that compute the automorphism group of a code, become expensive when the hull is
large, so the hull dimension measures their complexity \cite{sendrier1997}.
Codes with a trivial hull are the linear codes with complementary duals studied
by Massey \cite{massey1992}, and they can attain the Gilbert--Varshamov bound
\cite{sendrier2004}. In quantum coding the hull
dimension of a classical code determines the number of entangled pairs needed by
the corresponding entanglement-assisted quantum error-correcting code
\cite{brun2006,wildeBrun2008}, and the same mechanism is still used to obtain
quantum codes from constacyclic codes over rings \cite{zhang2026}.

The hulls of cyclic and negacyclic codes over finite fields were the first to be
described through the factorization of the modulus. Yang and Massey
characterized when a cyclic code has a complementary dual \cite{yangMassey1994},
Skersys computed the average dimension of the Euclidean hull over the cyclic
codes of a fixed length \cite{skersys2003}, and Sangwisut et al.\ determined the
Euclidean and Hermitian hull dimensions of cyclic and negacyclic codes in closed
form and counted the codes realizing each prescribed hull dimension
\cite{sangwisut2015}. Constacyclic codes, which contain both families, were
studied in the same manner: Jitman and Sangwisut obtained the average dimension
of the Hermitian hull of constacyclic codes over fields of square order
\cite{jitmanSangwisut2018}, and Debnath et al.\ extended the dimension formulas
to the $k$-Galois pairing for constacyclic codes over finite fields, together
with counts of the codes of a prescribed Galois hull dimension under
restrictions on the field parameters \cite{debnathConstacyclic2023}; a correction
appeared in \cite{debnathCorrection2023}. This line continues with the average
Galois hull dimension of constacyclic codes \cite{debnathAverage2025} and with
the small values a Galois hull dimension can attain \cite{debnathSmall2026}.

The Galois pairing also made the hull accessible for codes that are not
constacyclic. Liu and Pan characterized the Galois hull of an arbitrary linear
code over a finite field \cite{liuPan2020}, Cao constructed MDS codes with
Galois hulls of prescribed dimensions \cite{cao2021}, and Fu and Liu treated
Galois self-orthogonal constacyclic codes \cite{fuLiu2022}. More recently,
generator descriptions have been derived that avoid passing to the dual code:
generator polynomial matrices for the Galois hulls of multi-twisted codes
\cite{eldinSole2026}, and generator matrices for the intersection of a linear
code with its Galois dual \cite{eldinLeroy2026}, the operation that defines the
hull.

A separate question is not what the hull of one chosen code is, but how hull
dimensions are distributed over a family. For unrestricted linear codes over a
finite field, a mass formula gives the number of codes of a fixed length and
Euclidean hull dimension \cite{sendrier1997}, and the sequences of these counts
as the hull dimension grows have been studied recently \cite{bouyuklieva2026}.
Distributions and asymptotics have also been obtained for double cyclic codes
\cite{gao2023} and for double circulant and four circulant codes with small hull
dimension \cite{aliabadi2026}. The averages computed in \cite{skersys2003},
\cite{jitmanSangwisut2018} and \cite{debnathAverage2025} are the first moments of
such distributions and do not determine them.

When the alphabet is a ring rather than a field, the factorization theory
changes. The structure theory of cyclic and negacyclic codes over finite chain
rings \cite{dinhLopezPermouth2004} and the Hamming-distance theory over such
rings \cite{nortonSalagean2000} provide the background for the analyses below.
Talbi et al.\ studied the hulls of cyclic serial codes over a finite chain ring
\cite{talbi2022}; Jitman et al.\ treated cyclic codes over $\mathbb{Z}_4$, where
the length need not be coprime to the characteristic \cite{jitmanZ4}; and Pathak
and Sharma covered the oddly even case, which requires separate arguments
\cite{pathakZ4}. Constacyclic codes over a finite non-chain ring, with
applications to quantum codes, were considered in \cite{zhang2026}. Closest to
the present setting, Debnath et al.\ studied $k$-Galois hulls of constacyclic
codes over affine algebra rings, exhibiting generators of the duals and hulls, a
hull-dimension formula, LCD criteria and quantum-code applications
\cite{debnathAffine2026}.

Two kinds of results have therefore been obtained: descriptions of the hull of
an individual code, and counts or averages of a single statistic of a family.
Neither kind determines the joint distribution of the code dimension and the
hull dimension, and it is convenient to separate three questions: (i) which hull
dimensions occur; (ii) how many codes realize a prescribed hull dimension; and
(iii) how many codes realize each pair $(\kappa,\eta)$ formed by a code
dimension $\kappa$ and a hull dimension $\eta$. Questions (i) and (ii) are
answered for cyclic and negacyclic codes in \cite{sangwisut2015} and, over
finite fields, for constacyclic codes in \cite{debnathConstacyclic2023}, where
the counts require restrictions on the field parameters, while for unrestricted
linear codes with respect to the Euclidean dual, question (iii) is answered by
\cite{sendrier1997} and \cite{bouyuklieva2026}. Over an affine algebra only
question (i) has been answered, by the hull-dimension formula of
\cite{debnathAffine2026}. Such a formula does not record the code dimensions at
which the hull dimensions occur, and a hull-only count cannot be aggregated over
the different factor selections into a joint distribution. Answering question
(iii) in this setting therefore requires carrying the code dimension along with
the hull dimension and treating cycles of the factor permutation that are longer
than the fixed points and two-cycles that arise in the Euclidean and Hermitian
cases.

Motivated by this gap, we determine in this paper the joint distribution of code
dimension and $k$-Galois hull dimension for simple-root constacyclic codes over
a fixed square-free affine algebra $\A$ with labeled field components. The
counted objects are labeled factor selections, in bijection with the
$\lambda$-constacyclic codes over $\A$ (Lemmas~\ref{lem:complete}
and~\ref{lem:global-code}); each is a distinct submodule of the fixed ambient
space, and no isometry classes are formed. The output is a bivariate generating
polynomial in which the exponents of $u$ and $z$ record the code dimension
$\kappa$ and the $k$-Galois hull dimension $\eta_k$, so that its coefficients
answer the three questions above simultaneously. The method is to
decompose $\A$ into its field components, to translate $\lambda$-constacyclic
codes into labeled selections of irreducible factors of the component moduli,
and to follow the $k$-Galois dual through the Frobenius-twisted reciprocal
permutation of those factors. Both statistics are additive over the orbits of
this permutation, and this is what makes the enumeration factorize.

The results of the paper are the following.
(i) The $k$-Galois dual of a component code is generated by the normalized
inverse-Frobenius reciprocal of the parity-check polynomial and is constacyclic
with twist $\lambda_s^{-p^{d_s-k}}$ (Theorem~\ref{thm:dual-generator}), obtained
from the classical Euclidean description of a constacyclic dual
(Lemma~\ref{lem:ordinary-reciprocal}) by the coordinatewise inverse Frobenius
(Proposition~\ref{prop:dual-translation}).
(ii) Under the compatibility condition \eqref{eq:compat} the reciprocal
permutes the irreducible factors of the component modulus. The hull support of a
code is then a cyclic boundary statistic on the orbits of that permutation, and
the hull dimension is a weighted count of the resulting $1$-to-$0$ transitions
(Theorem~\ref{thm:hull-support} and Proposition~\ref{prop:boundary}).
(iii) Each orbit contributes a two-state weighted transfer matrix whose trace,
the orbit polynomial, also has a closed binomial form
(Propositions~\ref{prop:transfer-matrix}--\ref{prop:closed}).
(iv) The exact labeled joint $(\kappa,\eta)$ enumerator over the CRT product,
with orbit weights $w_O=m_sd_O$, is
\[
\Epoly(u,z)=\sum_{C}u^{\kappa(C)}z^{\eta_k(C)}
 =\prod_{s}\prod_{O\in\Os}\operatorname{tr}\bigl(T_{w_O}(u,z)^{a_O}\bigr)
\]
(Theorem~\ref{thm:central}), where the sum runs over all $\lambda$-constacyclic
codes over $\A$. Items (i)--(iii) assemble established descriptions of component
duals and closed-walk counting; the joint enumeration (iv) combines them into
one product formula. From it the two marginal distributions, the exact LCD
count, and the mean, variance and conditional distribution of the hull dimension
follow by specialization
(Corollaries~\ref{cor:total}--\ref{cor:conditional-variance}), while the
extension-degree weights $w_O$ convert component dimensions into dimensions over
the common base field.

The enumeration is exact: the coefficients of $\Epoly(u,z)$ are integers, the
conditional moments are reduced rationals, and for $q=32$, $n=31$ the orbit
recurrence evaluates all $2^{31}$ labeled factor selections without enumerating
them. For the families in characteristics $2$, $3$ and $5$ listed in
Section~\ref{sec:verification}, every $k$-Galois hull dimension computed
directly by finite-field linear algebra agrees with the orbit formula, and two
published benchmark counts \cite{sangwisut2015} are recovered. The orbit product
assumes simple-root moduli and compatible twists, and codes are counted as
labeled selections rather than isometry classes; the scope is stated at the end
of Section~\ref{sec:conclusion}. Figure~\ref{fig:pipeline} summarizes the
workflow. Section~\ref{sec:setting} fixes the algebraic setting and the
correspondence between constacyclic codes and labeled factor selections,
Section~\ref{sec:duality} develops component duality and the reciprocal-orbit
description of hull support, Section~\ref{sec:enumeration} derives the
enumerators with their statistical consequences, Section~\ref{sec:examples}
presents the examples and the computational validation, and
Section~\ref{sec:conclusion} collects the conclusions.
